\section{TODO Tasks}\label{sec:TODO}

This is an unordered list of tasks that still need to be done the Pollux Lean
formalization.

\begin{itemize}
\item \textbf{Reserved field numbers.} Protobuf doesn't recommend reusing field
  numbers for new fields (i.e.\ removing a field, then later adding a field with
  that same number), and that also poses a challenge to Pollux. If a field were
  added which uses and old field number, Pollux would assume that the intent is
  for the new field to be compatible with the old one. Since we expect the
  compatibility relation to be transitive, any evolving derivation of that
  message would attempt to force a connection between two fields that are
  semantically unrelated.

  Fortunately, Protobuf has a mechanism for this, reserved field numbers. You
  can set specific field numbers as reserved, then the protobuf compiler will
  ensure that you don't add a field with that number to the message in the
  future. Pollux will also enforce this, even though the Lean formalization
  doesn't track it \emph{yet}.

\item \textbf{Recursive Messages.} Protobuf messages (and thus descriptors) can
  be recursive. The current lean formalization enforces a tree structure on the
  descriptors. This is nice for recursion over the descriptors, but does mean we
  can't model all descriptors. The Protobuf data set does contain recursive and
  mutually recursive descriptors. It seems that these are more common in tools
  which analyze Protobuf descriptors, likely because \texttt{descriptor.proto},
  which models Protobuf descriptors as a Protobuf struct, is recursive. If we
  want to transmit protobuf descriptors into the Pollux lean formalization using
  \texttt{descriptor.proto} (a reasonable choice), we'll have to expand the
  descriptor representation to be a flat-map of named messages. See
  Table~\ref{tab:recur-summary} and Table~\ref{tab:recur-repos} for information
  about recursive messages in the Protobuf data set. This should be explored as
  a flat-map for descriptors (messages have to bottom-out, so this isn't a
  concern there) and see Section~\ref{sec:type-theory} for a discussion on avoid
  co-induction.

\item \textbf{Enums.} Protobuf supports enums, which shouldn't be hard to
  support in the future, but aren't present currently.

\item \textbf{Strings as UTF-8.} All valid Protobuf strings are encoded as UTF-8
  strings. Right now, the Lean formalization uses the built-in \texttt{String}
  type, and it has the ability to check if something is valid UTF-8 with
  \texttt{String.fromUTF8? (toUTF8 s) = some s}.

\item \textbf{Protobuf Serializer.} We need to implement a Protobuf serializer
  in Lean.

\item \textbf{Protobuf Parser.} We need to implement a Protobuf parser in Lean.

\end{itemize}

These tasks need to be accomplished on the theory side.

\begin{itemize}
\item \textbf{Enums.} Protobuf supports enums, which shouldn't be hard to
  support in the future, but aren't present currently.

\item \textbf{Recursive Messages.} We need to decide on how to handle these on
  the theory side.
\end{itemize}

% Local Variables:
% citar-bibliography: ("../pollux.bib")
% TeX-master: "../pollux.tex"
% TeX-command-default: "Make"
% jinx-local-words: "protobuf"
% End:
